% GENERATED FILE. Edit canonical lessons, then run npm run generate:books.
% canonical-source-hash: fnv1a64:db0d3becf0f6ac8b
% canonical-lessons: BN-C249-moshla, BN-C249-jire, BN-C249-dhone, BN-C249-pudina, BN-C249-kaju

\chapter{Spice, Cumin, Coriander, Mint, Cashew}
\label{ch:bn-spice-cumin-coriander-mint-cashew}
\hlchaptermodality{249}

\begin{chapteropening}
\textbf{By the end of this chapter:} \emph{I can say spice, cumin, coriander, mint and a cashew.}

Complete the last lesson of chapter 249: I can say spice, cumin, coriander, mint and a cashew.
\end{chapteropening}

\section[\texorpdfstring{\bn{মশলা} (môshlā)}{মশলা (môshlā)}]{\bn{মশলা} (môshlā) — spice}

\label{lesson:BN-C249-moshla}



\begin{quote}
Before the new one: say the Bengali for a parcel, then the Bengali for a magazine.
\end{quote}

\subsection*{You'll want to know: \bn{মশলা}}
\textbf{\bn{মশলা}} — \emph{môshlā} — \textquotedblleft{}spice\textquotedblright{}.

\begin{grammarlens}[title={What you've built}]
One more word for this chapter.
\end{grammarlens}

\subsection*{Your turn}
\begin{itemize}
\raggedright
  \item \emph{Say it:} \emph{môshlā}
  \item \emph{Say it:} \emph{môshlā}, once more
  \item \emph{Say it:} say \emph{môshlā}
  \item \emph{Recall:} say the Bengali for a parcel, then the Bengali for a magazine, then say \emph{môshlā} again
\end{itemize}

\begin{tcolorbox}[breakable,colback=teal!4,colframe=teal!35!black,title={Before you move on}]
What does \bn{মশলা} mean? (\textquotedblleft{}Spice\textquotedblright{}.) Say it once more.
\end{tcolorbox}
\section[\texorpdfstring{\bn{জিরে} (jire)}{জিরে (jire)}]{\bn{জিরে} (jire) — cumin}

\label{lesson:BN-C249-jire}



\begin{quote}
Before the new one: say the Bengali for a magazine, then the Bengali for spice.
\end{quote}

\subsection*{You'll want to know: \bn{জিরে}}
\textbf{\bn{জিরে}} — \emph{jire} — \textquotedblleft{}cumin\textquotedblright{}.

\begin{grammarlens}[title={What you've built}]
One more word for this chapter.
\end{grammarlens}

\subsection*{Your turn}
\begin{itemize}
\raggedright
  \item \emph{Say it:} \emph{jire}
  \item \emph{Say it:} \emph{jire}, once more
  \item \emph{Say it:} say \emph{jire}
  \item \emph{Recall:} say the Bengali for a magazine, then the Bengali for spice, then say \emph{jire} again
\end{itemize}

\begin{tcolorbox}[breakable,colback=teal!4,colframe=teal!35!black,title={Before you move on}]
What does \bn{জিরে} mean? (\textquotedblleft{}Cumin\textquotedblright{}.) Say it once more.
\end{tcolorbox}
\section[\texorpdfstring{\bn{ধনে} (dhône)}{ধনে (dhône)}]{\bn{ধনে} (dhône) — coriander}

\label{lesson:BN-C249-dhone}



\begin{quote}
Before the new one: say the Bengali for spice, then the Bengali for cumin.
\end{quote}

\subsection*{You'll want to know: \bn{ধনে}}
\textbf{\bn{ধনে}} — \emph{dhône} — \textquotedblleft{}coriander\textquotedblright{}.

\begin{grammarlens}[title={What you've built}]
One more word for this chapter.
\end{grammarlens}

\subsection*{Your turn}
\begin{itemize}
\raggedright
  \item \emph{Say it:} \emph{dhône}
  \item \emph{Say it:} \emph{dhône}, once more
  \item \emph{Say it:} say \emph{dhône}
  \item \emph{Recall:} say the Bengali for spice, then the Bengali for cumin, then say \emph{dhône} again
\end{itemize}

\begin{tcolorbox}[breakable,colback=teal!4,colframe=teal!35!black,title={Before you move on}]
What does \bn{ধনে} mean? (\textquotedblleft{}Coriander\textquotedblright{}.) Say it once more.
\end{tcolorbox}
\section[\texorpdfstring{\bn{পুদিনা} (pudinā)}{পুদিনা (pudinā)}]{\bn{পুদিনা} (pudinā) — mint}

\label{lesson:BN-C249-pudina}



\begin{quote}
Before the new one: say the Bengali for cumin, then the Bengali for coriander.
\end{quote}

\subsection*{You'll want to know: \bn{পুদিনা}}
\textbf{\bn{পুদিনা}} — \emph{pudinā} — \textquotedblleft{}mint\textquotedblright{}.

\begin{grammarlens}[title={What you've built}]
One more word for this chapter.
\end{grammarlens}

\subsection*{Your turn}
\begin{itemize}
\raggedright
  \item \emph{Say it:} \emph{pudinā}
  \item \emph{Say it:} \emph{pudinā}, once more
  \item \emph{Say it:} say \emph{pudinā}
  \item \emph{Recall:} say the Bengali for cumin, then the Bengali for coriander, then say \emph{pudinā} again
\end{itemize}

\begin{tcolorbox}[breakable,colback=teal!4,colframe=teal!35!black,title={Before you move on}]
What does \bn{পুদিনা} mean? (\textquotedblleft{}Mint\textquotedblright{}.) Say it once more.
\end{tcolorbox}
\section[\texorpdfstring{\bn{কাজু} (kāju)}{কাজু (kāju)}]{\bn{কাজু} (kāju) — a cashew}

\label{lesson:BN-C249-kaju}



\begin{quote}
Before the new one: say the Bengali for coriander, then the Bengali for mint.
\end{quote}

\subsection*{You'll want to know: \bn{কাজু}}
\textbf{\bn{কাজু}} — \emph{kāju} — \textquotedblleft{}a cashew\textquotedblright{}.

\begin{grammarlens}[title={What you've built}]
One more word for this chapter.
\end{grammarlens}

\subsection*{Your turn}
\begin{itemize}
\raggedright
  \item \emph{Say it:} \emph{kāju}
  \item \emph{Say it:} \emph{kāju}, once more
  \item \emph{Say it:} say \emph{kāju}
  \item \emph{Recall:} say the Bengali for coriander, then the Bengali for mint, then say \emph{kāju} again
\end{itemize}

\begin{tcolorbox}[breakable,colback=teal!4,colframe=teal!35!black,title={Before you move on}]
What does \bn{কাজু} mean? (\textquotedblleft{}A cashew\textquotedblright{}.) Say it once more.
\end{tcolorbox}
